% !PRINTED 182
\weberpage{182}
%<<CONT
\weq{1}{\begin{array}{l}
\sigma(v,\,2K,\,2iK') = \displaystyle\sum_{0,\infty}^{\nu} A_\nu\,\dfrac{v^{2\nu+1}}{(2\nu+1)!},\\[8pt]
\sigma_{00}(v,\,2K,\,2iK') = \displaystyle\sum_{0,\infty}^{\nu} B_\nu\,\dfrac{v^{2\nu}}{(2\nu)!},\\[8pt]
\sigma_{01}(v,\,2K,\,2iK') = \displaystyle\sum_{0,\infty}^{\nu} C_\nu\,\dfrac{v^{2\nu}}{(2\nu)!},\\[8pt]
\sigma_{10}(v,\,2K,\,2iK') = \displaystyle\sum_{0,\infty}^{\nu} D_\nu\,\dfrac{v^{2\nu}}{(2\nu)!},
\end{array}}
und die $A_\nu$, $B_\nu$, $C_\nu$, $D_\nu$ sind die zu bestimmenden Funktionen von $\varkappa^2$ oder von $\omega$. Fassen wir sie zunächst als Funktionen von $\omega$ auf, so ergeben die linearen Transformationen [\S\,45, (3)]:
\weq{2}{\begin{array}{l}
A_\nu\left(-\dfrac{1}{\omega}\right) = (-1)^\nu A_\nu(\omega),\\[6pt]
B_\nu\left(-\dfrac{1}{\omega}\right) = (-1)^\nu B_\nu(\omega),\\[6pt]
C_\nu\left(-\dfrac{1}{\omega}\right) = (-1)^\nu D_\nu(\omega),\\[6pt]
D_\nu\left(-\dfrac{1}{\omega}\right) = (-1)^\nu C_\nu(\omega).
\end{array}}
\weq{3}{\begin{array}{l}
\varkappa'^{2\nu} A_\nu(\omega+1) = A_\nu(\omega),\\[4pt]
\varkappa'^{2\nu} B_\nu(\omega+1) = C_\nu(\omega),\\[4pt]
\varkappa'^{2\nu} C_\nu(\omega+1) = B_\nu(\omega),\\[4pt]
\varkappa'^{2\nu} D_\nu(\omega+1) = D_\nu(\omega).
\end{array}}
\weq{4}{\begin{array}{l}
A_\nu(\omega+2) = A_\nu(\omega),\\[4pt]
B_\nu(\omega+2) = B_\nu(\omega),\\[4pt]
C_\nu(\omega+2) = C_\nu(\omega),\\[4pt]
D_\nu(\omega+2) = D_\nu(\omega).
\end{array}}
Hieraus schließen wir auf die charakteristischen Eigenschaften dieser Koeffizienten als Funktionen von $\varkappa^2$. Zunächst erhellt aus der Bedeutung der Koeffizienten, daß für jeden Wert von $\varkappa^2$ mit etwaiger Ausnahme von $\varkappa^2 = 0, 1, \infty$ die Koeffizienten $A_\nu(\varkappa^2)$, $B_\nu(\varkappa^2)$, $C_\nu(\varkappa^2)$, $D_\nu(\varkappa^2)$ endliche und stetige Funktionen von $\varkappa^2$ sind.

Aber auch für $\varkappa^2 = 0$ bleiben diese Funktionen endlich und man kann ihre Werte leicht finden, wenn man in den Darstellungen des \S\,37 $q$ und mithin $\varkappa^2$ in Null übergehen läßt.
